\section{From a Prescribed Force to Terminal Continuation}
\label{f:sec:forced-argument-overview}
\label{f:guide:how-proof-works}

Fix one divergence-free datum \(u_0\), one viscosity \(\nu>0\), and one
prescribed smooth force \(f\).  Let \((u,p)\) be the corresponding maximal
classical solution on \([0,T_*)\).  The proof never changes this history.
Every derivative, finite rectangle, endpoint value, and restarted solution
below belongs to the trajectory generated by the original triple
\((u_0,\nu,f)\).

The force enters the proof in two linked forms.  Its solenoidal part
\(g=\Pp f\) acts in the projected velocity equation, while its gradient part
is fixed by the normalized pressure:
\[
 \partial_tu+\Pp\nabla\!\cdot(u\otimes u)-\nu\Delta u=g,
 \qquad
 p=R_iR_j(u_i u_j)-(-\Delta)^{-1}\nabla\!\cdot f.
\]
The source enters the velocity and pressure at the initial
construction and remains present in every differentiated row.

The Navier--Stokes scaling acts on the complete forced equation.  If
\((u,p,f)\) is a solution, then
\[
 u_\lambda(t,x)=\lambda u(\lambda^2t,\lambda x),\qquad
 p_\lambda(t,x)=\lambda^2p(\lambda^2t,\lambda x),\qquad
 f_\lambda(t,x)=\lambda^3f(\lambda^2t,\lambda x)
\]
solves the equation with transformed datum
\(u_{0,\lambda}(x)=\lambda u_0(\lambda x)\).  For every real \(s\),
\[
 \|\Lambda^su_\lambda(t)\|_2^2
 =\lambda^{2s-1}\|\Lambda^su(\lambda^2t)\|_2^2.
\]
The invariant spatial height is \(s=\tfrac12\).  This identifies
the critical row of the same forced differential structure; it does not
replace the prescribed history by its rescaling.

Write
\[
 V_n=\partial_t^nu,\qquad Q_n=\partial_t^np,
 \qquad F_n=\partial_t^nf,
\]
and let
\[
 B_n=\sum_{a=0}^n\binom na(V_a\cdot\nabla)V_{n-a}.
\]
Repeated differentiation of the original forced equation gives the complete
temporal row
\begin{align}
 V_{n+1}+B_n+\nabla Q_n&=\nu\Delta V_n+F_n,
 \label{f:eq:guide-forced-row}\\
 -\Delta Q_n
 &=\partial_i\partial_j\sum_{a=0}^n\binom na
      (V_a)_i(V_{n-a})_j-\nabla\!\cdot F_n.
 \label{f:eq:guide-forced-pressure}
\end{align}
Spatial differentiation distributes through both velocity factors and through
the prescribed source.  The fields
\[
 \partial_x^\alpha V_n,\qquad
 \partial_x^\alpha Q_n,\qquad
 \partial_x^\alpha F_n
\]
form the pressure-complete forced mixed jet.  They are coordinates of one
solution, rather than a collection of independently chosen functions.

\phantomsection\label{f:guide:intersection}
At a proposed finite maximal time \(T_*\), the finite object through temporal
depth \(N\) and spatial height \(M\) is the record
\[
 \cR_{N,M}[u_0,\nu,f;T_*]
 :=\bigl((V_n)_{n=0}^{N},(V_{n+1})_{n=0}^{N};
         (Q_n)_{n=0}^{N};(F_n)_{n=0}^{N}\bigr).
\]
Its scalar velocity norm on a strict horizon is
\[
 \mathfrak R_{N,M}(T)
 :=\sum_{n=0}^{N}\left(
 \|V_n\|_{L^2(0,T;H^{M+1})}^2
 +\|V_{n+1}\|_{L^2(0,T;H^{M-1})}^2\right).
\]
The datum construction places the record on \((0,T_*)\) and gives
\[
 \sup_{0<T<T_*}\mathfrak R_{N,M}(T)<\infty
 \qquad(N,M<\infty).
 \tag{A}
\]
Deleting outer temporal or spatial coordinates recovers the corresponding
smaller record.  As \(N\) and \(M\) range over all finite values, these blocks
exhaust every mixed derivative of the one forced history.  Restriction in time
gives the strict-horizon blocks, while monotone exhaustion identifies each
finite norm on \((0,T_*)\).

Letting the temporal and spatial indices range over all finite values gives
\[
 u\in\bigcap_{r,m\ge0}H^r(I;H^m_\sigma(\R^3)),
 \qquad I=(0,T_*).
 \tag{A$'$}
\]
Independently, the completed critical heights and the viscous diagonal relation
are derived from the same mixed jet.  They retain transport, pressure,
viscosity, and force work while exposing every finite spatial height.
Evaluating those identities on the datum-generated blocks and intersecting the
heights gives the complete spatial Sobolev column; the forced recurrence then
recovers the same intersection (A$'$).
There is no second family of fields: both constructions use the same
datum-generated block \(\cR_{N,M}[u_0,\nu,f;T_*]\) at every finite pair
\((N,M)\).  The independently derived critical-height relationship rectangle
is the finite system of coupled identities evaluated on those same
coordinates.

The viscous relation is visible already in one row.  Project
\eqref{f:eq:guide-forced-row}, apply \(\Lambda^s\), and use
\(\Delta=-\Lambda^2\):
\[
 \Lambda^sV_{n+1}+\nu\Lambda^{s+2}V_n
 =\Lambda^s\Pp(F_n-B_n).
 \tag{B}
\]
Squaring this equality does not separate the nonlinearity from the fluid.
The cross term is the derivative of the adjacent spatial norm, so
\begin{align}
 &\|\Lambda^sV_{n+1}\|_2^2
 +\nu^2\|\Lambda^{s+2}V_n\|_2^2
 +\nu\frac{d}{dt}\|\Lambda^{s+1}V_n\|_2^2
 \notag\\
 &\hspace{4cm}
 =\|\Lambda^s\Pp(F_n-B_n)\|_2^2.
 \tag{C}
\end{align}
The complementary projection gives the pressure allocation
\[
 \nabla Q_n=(I-\Pp)(F_n-B_n),
\]
and orthogonality recombines the complete source:
\[
 \|\Lambda^s\Pp(F_n-B_n)\|_2^2
 +\|\Lambda^s\nabla Q_n\|_2^2
 =\|\Lambda^s(F_n-B_n)\|_2^2.
 \tag{D}
\]
Thus the rectangle keeps viscosity, transport, pressure, and forcing in the
same estimate.

The critical-height tower is another reading of this viscous tether.  For
\[
 E_{r,s}=\frac12\|\Lambda^sV_r\|_2^2,
\]
the complete row is
\[
 \partial_tE_{r,s}+2\nu E_{r,s+1}
 =\mathcal P_{r,s}+\mathcal W_{r,s},
 \qquad
 \mathcal W_{r,s}
 =\operatorname{Re}\langle\Lambda^sV_r,\Lambda^sF_r\rangle.
 \tag{E}
\]
At \(r=0\) and \(s=\tfrac12\), differentiating (E) repeatedly exposes
\(E_{3/2},E_{5/2},\ldots\), while the differentiated transport, pressure,
and force work remains in the same completed row.  The tower then
describes the geometry among mixed-jet coordinates; it does not replace the
velocity field by a scalar height.
The first two spatial heights are, explicitly,
\begin{align*}
 H'+\nu\norm{\Lambda^{3/2}u}_2^2
 &=\mathcal P_{1/2}+\mathcal W_{1/2},\\
 E_{3/2}'+\nu\norm{\Lambda^{5/2}u}_2^2
 &=\mathcal P_{3/2}+\mathcal W_{3/2},
\end{align*}
and hence
\[
 H''=4\nu^2E_{5/2}
 -2\nu(\mathcal P_{3/2}+\mathcal W_{3/2})
 +\partial_t(\mathcal P_{1/2}+\mathcal W_{1/2}).
\]
These are the actual \(s=\tfrac12\) and \(s=\tfrac32\) rows.  Their
transport, pressure, viscosity, and force coordinates remain joined.

\phantomsection\label{f:guide:mixed-jet-map}
For finite \(N,M\), use the corresponding datum-generated block
\(\cR_{N,M}[u_0,\nu,f;T_*]\).
At the spatial orders required by (B)--(E), its restriction to a strict
horizon satisfies
\[
 \mathfrak R_{N,M}(T)
 \le \mathfrak R_{N,M}^*
 :=\sup_{0<S<T_*}\mathfrak R_{N,M}(S)<\infty.
 \tag{F}
\]
Restriction compatibility identifies every common row in two such
rectangles, while monotone exhaustion gives
\[
 \lim_{T\uparrow T_*}\mathfrak R_{N,M}(T)
 =\mathfrak R_{N,M}^*<\infty.
 \tag{G}
\]
The bound may depend on the selected rectangle.  No estimate uniform in the
infinitely many indices is used.  The bounds in (G) are scalar norms of the
finite blocks whose mixed-index exhaustion gives (A$'$).  The
independently derived identities in (E), evaluated on those same coordinates,
give the completed heights.  Exhausting their spatial heights and applying the
Fourier Sobolev split gives the spatial column, from which
\eqref{f:eq:guide-forced-row}--\eqref{f:eq:guide-forced-pressure} recover
(A$'$) again.  The positive-sign square (C), together with (D), relates the
adjacent temporal, viscous, and pressure faces.  Sobolev
multiplication is applied within that same finite rectangle to its binomial
transport face \(B_n\); it does not replace the nonlinear term or introduce a
second premise.

The mixed-index exhaustion and the critical-height relationships use the same
finite blocks.  The former records every mixed derivative; the
latter recover every spatial row and return through the forced recurrence.
Pairing, squaring, and integrating (B)--(E) inside any selected rectangle gives
its complete-equation estimates, with the nonlinearity retained as one face of
that forced rectangle.

\phantomsection\label{f:guide:finite-estimates}
The datum rectangles now carry the same complete equation to the terminal
step.  Their zeroth temporal faces give, for every finite \(m\),
\[
 u\in L^2(0,T_*;H^{m+2}).
\]
Sobolev multiplication and the prescribed force then give
\begin{align*}
 \|u_t\|_{L^1(0,T_*;H^m)}
 &\le \nu\sqrt{T_*}\,\|u\|_{L^2(0,T_*;H^{m+2})}\\
 &\quad+C_m\|u\|_{L^2(0,T_*;H^{m+2})}^2
 +\|g\|_{L^1(0,T_*;H^m)}<\infty.
\end{align*}
Hence
\[
 u\in W^{1,1}(0,T_*;H^m)
 \hookrightarrow C([0,T_*];H^m)
 \qquad(m\ge0).
\]
All these traces come from the same velocity and determine one
terminal state
\[
 u_*=u(T_*)\in\bigcap_{m\ge0}H^m_\sigma(\R^3).
\]

Passing \eqref{f:eq:guide-forced-row}--\eqref{f:eq:guide-forced-pressure} to
\(T_*\) determines every temporal and pressure coordinate from \(u_*\) and
the values \(F_n(T_*)\).  Local theory from \(u_*\), driven by the same
prescribed force \(f(t)\), continues the same history past \(T_*\).
This contradicts maximality, so \(T_*=\infty\).

\phantomsection\label{f:guide:how-proof-ends}
